% =====================================================================
%  BÖLÜM 1 — Ölçüm Teorisinin Temelleri
%  Kaynaklar: Billingsley Böl.2-3; Klenke Böl.1-6; Shiryaev II.§1-§3
% =====================================================================

\chapter{Ölçüm Teorisinin Temelleri}

\bolumalinti{``The theory of measure is the foundation on which
  the modern theory of probability rests.''}{Andrey Kolmogorov (1903--1987)}

\noindent\textbf{Ön Koşullar:} Bölüm~0 (küme cebiri, sayılabilirlik, $\limsup/\liminf$).

\noindent\textbf{Bu Bölüm Sonraki Bölümlerde Kullanılır:}
Bölüm~2 (olasılık ölçüsü),
Bölüm~3 (rasgele değişkenler = ölçülebilir fonksiyonlar),
Bölüm~4 (beklenti = Lebesgue integrali).

\noindent\textbf{Tahmini Okuma Süresi:} $\approx 10$~saat.

\vspace{0.5\baselineskip}

\begin{kazanimlar}
Bu bölümün sonunda öğrenci:
\begin{itemize}
  \item $\sigma$-cebiri tanımını verebilmeli, örnekler üretebilmeli ve
        verilen bir aile tarafından üretilen $\sigma$-cebirini betimleyebilmeli,
  \item Borel $\sigma$-cebirini $\mathcal{B}(\mathbb{R})$ tanımlayabilmeli ve
        açık kümelerin, kapalı kümelerin, $F_\sigma$ ile $G_\delta$ kümelerinin
        Borel olduğunu gösterebilmeli,
  \item Ölçü aksiyomlarını ifade edebilmeli; monotonluk, alt-toplamsallık ve
        aşağıdan/yukarıdan süreklilik özelliklerini ispatlayabilmeli,
  \item Carathéodory genişleme teoremini ifade edebilmeli ve Lebesgue ölçüsünün
        bu teoremden nasıl elde edildiğini açıklayabilmeli,
  \item Ölçülebilir fonksiyon kavramını tanımlayabilmeli; basit fonksiyon
        yaklaşımını ve Lebesgue integralinin yapısını açıklayabilmeli,
  \item Monoton Yakınsama Teoremi, Fatou Lemması ve Dominante Yakınsama
        Teoremini ispatlayabilmeli ve bunları integral hesaplarında uygulayabilmeli,
  \item Fubini-Tonelli teoremini ifade edebilmeli ve çift integrallerin sırasını
        değiştirme koşulunu açıklayabilmeli.
\end{itemize}
\end{kazanimlar}

\begin{motivasyon}
Olasılık teorisinin Kolmogorov tarafından 1933'te aksiyomatik biçimde
kurulmadan önce, ``olasılık'' kavramı kesin matematiksel temellere sahip
değildi. Hangi olaylara olasılık atanabilir? Sonsuz çok olayı kapsayan
deneyler nasıl ele alınır? Lebesgue integralinin beklentiyle ilişkisi nedir?

Tüm bu soruların yanıtı ölçüm teorisinde yatar. Bu bölümde önce ``olası
olayların ailesinin'' hangi yapıya sahip olması gerektiğini ($\sigma$-cebiri),
ardından bu yapıya sayı atayan fonksiyonun özelliklerini (ölçü) inceleyeceğiz.

\medskip

\textbf{Köprü kurma.} Riemann integralini biliyorsunuz: fonksiyonu parçalara böler,
dikdörtgen alanları toplarsınız. Lebesgue integrali tam tersi yönde çalışır:
fonksiyonun değer aralığını böler, her değer kümesinin ``ölçüsünü'' hesaplar.
Bu fark küçük görünse de, Lebesgue integrali sınır alma ile değiştirilebilirlik
gibi son derece güçlü teoremler sağlar; Riemann integrali bunu sağlayamaz.
\end{motivasyon}

\begin{tarihselnot}
Henri Lebesgue, 1901 yılında doktora tezinde yeni bir integral teorisi kurdu.
Émile Borel daha önce belirli küme ailelerine ``ölçü'' atamayı düşünmüştü;
Lebesgue bu fikri sistematikleştirdi. Carathéodory ise 1914'te dış ölçü
yoluyla genişleme teoremini formüle etti. Kolmogorov, 1933'te tüm bu yapıyı
olasılık teorisine uyguladı: olasılık uzayı $(\Omega, \mathcal{F}, \mathbb{P})$
modern matematiksel olasılığın temel nesnesi oldu.
\end{tarihselnot}


% =====================================================================
\section{$\sigma$-Cebirleri ve Borel Kümeler}
% =====================================================================

\begin{kesif}[$\sigma$-cebiri neden gerekli?]
Bir madeni para \emph{sonsuz kez} atılıyor. ``İlk tura $n$-inci atışta gelir''
olayını $A_n$ ile gösterin. ``En az bir kez tura gelir'' olayı nedir?
Cevabınızı aşağıya yazın.
\kesifcevap{``En az bir kez tura gelir'' = $\bigcup_{n=1}^{\infty} A_n$.
  Bu sonsuz birleşimi içeren her aile $\sigma$-cebiri olmak zorundadır;
  aksi takdirde bu olaya olasılık atamak tutarsızlığa yol açar.
  $\sigma$-cebirinin kapalılık koşulu buradan doğrudan çıkar.}
Olasılık tanımlamak istediğimiz olaylar ailesi bu tür işlemler altında kapalı
olmalıdır. Bu koşulun kesin matematiksel karşılığı $\sigma$-cebiriyidir.
\end{kesif}

\begin{tanimgerekce}
Neden sadece sonlu işlemler değil, \emph{sayılabilir} işlemler istiyoruz?
Çünkü sonsuz deneyler (bir madeni paranın sonsuz kez atılması, Poisson süreci)
doğal olarak sayılabilir sonsuz birleşimler üretir. Sayılamaz birleşimlere
izin versek, her tekil kümenin birleşimi tüm $\mathbb{R}$'yi verir; bu kümeler
üzerinde Lebesgue ölçüsü gibi makul bir ölçü tanımlamak imkânsız olur.
\end{tanimgerekce}

\begin{tanim}[$\sigma$-Cebiri]\label{def:sigmacebiri}
$\Omega$ boşalmayan bir küme olsun. $\mathcal{F} \subseteq \mathcal{P}(\Omega)$
alt kümeler ailesine \textbf{$\sigma$-cebiri} \emph{($\sigma$-algebra,
$\sigma$-field)} denir; eğer:
\begin{enumerate}[label=\textbf{(\Alph*)}]
  \item $\Omega \in \mathcal{F}$,
  \item $A \in \mathcal{F} \Rightarrow A^c \in \mathcal{F}$ \quad(tümleme altında kapalı),
  \item $A_1, A_2, \ldots \in \mathcal{F} \Rightarrow \bigcup_{n=1}^\infty A_n \in \mathcal{F}$
        \quad(sayılabilir birleşim altında kapalı).
\end{enumerate}
$(\Omega, \mathcal{F})$ çiftine \textbf{ölçülebilir uzay} \emph{(measurable space)} denir.
\end{tanim}

\begin{sonuc}[Türetilen Özellikler]\label{cor:sigmaturev}
$\mathcal{F}$ bir $\sigma$-cebiri ise:
\begin{enumerate}[label=(\roman*)]
  \item $\emptyset \in \mathcal{F}$,
  \item $A_1, A_2, \ldots \in \mathcal{F} \Rightarrow \bigcap_{n=1}^\infty A_n \in \mathcal{F}$,
  \item $A, B \in \mathcal{F} \Rightarrow A \setminus B \in \mathcal{F}$,
  \item $A, B \in \mathcal{F} \Rightarrow A \triangle B \in \mathcal{F}$.
\end{enumerate}
\end{sonuc}

\begin{proof}
(i) $\Omega \in \mathcal{F}$'den (A), tümlemeyle $\emptyset = \Omega^c \in \mathcal{F}$.
(ii) De~Morgan: $\bigcap_n A_n = \left(\bigcup_n A_n^c\right)^c$;
her $A_n^c \in \mathcal{F}$ (B'den), birleşim $\mathcal{F}$'de (C), tümleyen $\mathcal{F}$'de (B).
(iii) $A \setminus B = A \cap B^c$; (ii)'nin özel hali.
(iv) $(A \setminus B) \cup (B \setminus A)$; (iii) ve (C)'den.
\end{proof}

\begin{ornek}[$\sigma$-cebiri örnekleri]\label{ex:sigmaoernekler}
\begin{enumerate}[label=(\alph*)]
  \item \textbf{Trivial (önemsiz) $\sigma$-cebiri:} $\mathcal{F} = \{\emptyset, \Omega\}$.
        Her $\Omega$ üzerinde en küçük $\sigma$-cebiriyidir.
  \item \textbf{Kuvvet kümesi:} $\mathcal{F} = \mathcal{P}(\Omega)$.
        En büyük $\sigma$-cebiriyidir.
  \item $\Omega = \{1,2,3,4\}$, $\mathcal{F} = \{\emptyset, \{1,2\}, \{3,4\}, \Omega\}$.
        Doğrulayın: (A), (B), (C) sağlanıyor.
  \item $\mathcal{F} = \{A \subseteq \Omega : A \text{ sayılabilir ya da } A^c \text{ sayılabilir}\}$.
        $\Omega$'nın herhangi bir güçte olabileceği bu ``sayılabilir-kosayılabilir''
        $\sigma$-cebiri standart bir örnektir.
\end{enumerate}
\end{ornek}

\begin{kritiksorgu}[Sonlu aileler yeterli değil mi?]
Tanımda sayılabilir birleşimi sonlu birleşimle değiştirseydik, elde edeceğimiz
yapıya \textbf{cebir} \emph{(algebra)} denir. Bir cebir, $\sigma$-cebiri olmak
zorunda değildir: $\mathcal{F}_0 = \{A \subseteq [0,1] : A \text{ ya da } A^c
\text{ sonlu birleşim aralıktır}\}$ bir cebirdir ama $\sigma$-cebiri değildir.
\end{kritiksorgu}

\subsection{Üretilmiş $\sigma$-Cebiri}

\begin{tanim}[Üretilmiş $\sigma$-Cebiri]\label{def:sigmauretilensi}
$\mathcal{E} \subseteq \mathcal{P}(\Omega)$ bir aile olsun.
$\mathcal{E}$'yi içeren tüm $\sigma$-cebirlerin kesişimine
$\mathcal{E}$ tarafından \textbf{üretilen $\sigma$-cebiri} denir ve
$\sigma(\mathcal{E})$ ile gösterilir:
\[
  \sigma(\mathcal{E}) \coloneqq \bigcap \{\mathcal{F} : \mathcal{F} \text{ bir } \sigma\text{-cebiri}, \ \mathcal{E} \subseteq \mathcal{F}\}.
\]
\end{tanim}

\begin{onerme}[$\sigma(\mathcal{E})$'nin iyi tanımlılığı]\label{prop:sigmaiyitanim}
$\sigma(\mathcal{E})$:
\begin{enumerate}[label=(\roman*)]
  \item Bir $\sigma$-cebiriyidir.
  \item $\mathcal{E}$'yi içeren en küçük $\sigma$-cebiriyidir:
        $\mathcal{E} \subseteq \mathcal{G}$ ve $\mathcal{G}$ $\sigma$-cebiri ise
        $\sigma(\mathcal{E}) \subseteq \mathcal{G}$.
\end{enumerate}
\end{onerme}

\begin{proof}
(i) $\sigma$-cebiri ailesinin keyfi kesişimi $\sigma$-cebiriyidir:
(A), (B), (C) aksiyomlarının her biri kesişimde korunur.
(ii) $\mathcal{G}$ kesişimde yer alan ailelerden biridir; dolayısıyla
kesişim $\subseteq \mathcal{G}$.
\end{proof}

\subsection{Borel $\sigma$-Cebiri}

\begin{tanim}[Borel $\sigma$-Cebiri]\label{def:borel}
$\mathbb{R}$ üzerindeki \textbf{Borel $\sigma$-cebiri}
\[
  \mathcal{B}(\mathbb{R}) \coloneqq \sigma\bigl(\{(a,b) : a < b,\ a,b \in \mathbb{R}\}\bigr)
\]
açık aralıklar ailesi tarafından üretilen $\sigma$-cebiriyidir.
$\mathcal{B}(\mathbb{R})$'nin elemanlarına \textbf{Borel kümesi} denir.
\end{tanim}

\begin{teorem}[Borel $\sigma$-Cebirinin Eşdeğer Üreticileri]\label{thm:boreldenk}
$\mathcal{B}(\mathbb{R})$ aşağıdakilerin her biri tarafından üretilir:
\begin{enumerate}[label=(\roman*)]
  \item Açık aralıklar: $\{(a,b) : a<b\}$
  \item Kapalı aralıklar: $\{[a,b] : a<b\}$
  \item Yarı-açık aralıklar: $\{(a,b] : a<b\}$\ veya\ $\{[a,b) : a<b\}$
  \item $\mathbb{R}$'nin açık kümeleri
  \item $(-\infty, b]$ biçimindeki kümeler ($b \in \mathbb{R}$)
\end{enumerate}
\end{teorem}

\begin{proof}
(v) için göstereceğiz; diğerleri benzerdir.

$\mathcal{E} = \{(-\infty, b] : b \in \mathbb{R}\}$ olsun. Önce $\mathcal{E} \subseteq \mathcal{B}(\mathbb{R})$:
\[
  (-\infty, b] = \bigcap_{n=1}^\infty (-\infty, b + 1/n)
  = \bigcap_{n=1}^\infty \Bigl[(-n, b+1/n)\Bigr]^c \cup \ldots
\]
Daha doğrudan: $(-\infty, b] = \mathbb{R} \setminus (b, \infty)$ ve
$(b,\infty) = \bigcup_{n=1}^\infty (b, b+n)$ açık aralıkların birleşimi;
dolayısıyla Borel. Tümleni de Borel.

Tersine, $(a,b) = (-\infty, b) \setminus (-\infty, a]$ ve
$(-\infty, b) = \bigcup_{n=1}^\infty (-\infty, b - 1/n]$; dolayısıyla
$\sigma(\mathcal{E})$ açık aralıkları içerir, yani $\mathcal{B}(\mathbb{R}) \subseteq \sigma(\mathcal{E})$.
İki taraflı içerme ile $\sigma(\mathcal{E}) = \mathcal{B}(\mathbb{R})$.
\end{proof}

\begin{notkutu}
Borel kümeleri çok geniş bir ailedir. Açık, kapalı, $G_\delta$ (açıkların sayılabilir
kesişimi), $F_\sigma$ (kapalıların sayılabilir birleşimi), $G_{\delta\sigma}$, \ldots
hepsi Boreldir. Buna karşın Borel olmayan Lebesgue-ölçülebilir kümeler
mevcuttur (Lebesgue $\sigma$-cebirinin Borel $\sigma$-cebirinden geniş olduğunun kanıtı).
\end{notkutu}

\begin{disiplinlerarasibaglan}[Topoloji ve Fonksiyonel Analiz]
$\mathcal{B}(\mathbb{R})$, metrik uzaylarda sürekli fonksiyonların doğal
tanım alanıdır. Bir $f: \mathbb{R} \to \mathbb{R}$ fonksiyonu sürekli ise,
her açık kümenin ters görüntüsü açıktır; dolayısıyla her Borel kümesinin
ters görüntüsü Boreldir. Bu gözlem Bölüm~3'teki ölçülebilir fonksiyon
kavramının temel örneğini oluşturur.
\end{disiplinlerarasibaglan}

\begin{mikroozet}[§1.1 özeti]
$\sigma$-cebiri: tümleme ve sayılabilir birleşim altında kapalı küme ailesi.
$\sigma(\mathcal{E})$: $\mathcal{E}$'yi içeren en küçük $\sigma$-cebiri.
$\mathcal{B}(\mathbb{R})$: açık aralıkların ya da açık kümelerin ürettiği $\sigma$-cebiri.
\end{mikroozet}


% =====================================================================
\section{Ölçüler ve Olasılık Ölçüleri}
% =====================================================================

\begin{tanimgerekce}
Bir $\sigma$-cebirine sayı atayan fonksiyon neden $\sigma$-toplamsallık
gerektiriyor? Çünkü ``uzunluk'', ``alan'', ``hacim'' ve ``olasılık'' gibi
kavramların hepsinde \emph{ayrık parçaların toplamı bütüne eşittir.} Sonlu
toplamsallık tek başına yeterli değil: sayılabilir sonsuz birleşimlerde
süreklilik özelliği istiyoruz.
\end{tanimgerekce}

\begin{tanim}[Ölçü]\label{def:olcu}
$(\Omega, \mathcal{F})$ bir ölçülebilir uzay ve
$\mu: \mathcal{F} \to [0,+\infty]$ bir fonksiyon olsun.
$\mu$'ya \textbf{ölçü} \emph{(measure)} denir; eğer:
\begin{enumerate}[label=\textbf{(\Alph*)}]
  \item $\mu(\emptyset) = 0$,
  \item \textbf{$\sigma$-toplamsallık} \emph{($\sigma$-additivity)}: $(A_n)_{n \geq 1}$
        ikili ayrık ve $A_n \in \mathcal{F}$ ise
        \[
          \mu\!\left(\bigcup_{n=1}^\infty A_n\right) = \sum_{n=1}^\infty \mu(A_n).
        \]
\end{enumerate}
$(\Omega, \mathcal{F}, \mu)$ üçlüsüne \textbf{ölçüm uzayı} \emph{(measure space)} denir.
\end{tanim}

\begin{tanim}[Olasılık Ölçüsü]\label{def:olasilikolucu}
$\mathbb{P}: \mathcal{F} \to [0,1]$ bir ölçüdür ve $\mathbb{P}(\Omega) = 1$
ise $\mathbb{P}$'ye \textbf{olasılık ölçüsü} \emph{(probability measure)} denir.
$(\Omega, \mathcal{F}, \mathbb{P})$ üçlüsüne \textbf{olasılık uzayı} denir.
\end{tanim}

\begin{notkutu}
Olasılık ölçüsü bir ölçünün özel halidir: değerleri $[0,1]$'e kısıtlı ve
tüm uzayın ölçüsü $1$'dir. Bölüm~2'de olasılık uzayını ayrıntılı inceleyeceğiz.
\end{notkutu}

\begin{teorem}[Ölçünün Temel Özellikleri]\label{thm:olcuozellik}
$(\Omega, \mathcal{F}, \mu)$ bir ölçüm uzayı olsun. O zaman:
\begin{enumerate}[label=(\roman*)]
  \item \textbf{Sonlu toplamsallık:} $A_1, \ldots, A_n$ ikili ayrık ise
        $\mu\!\left(\bigcup_{k=1}^n A_k\right) = \sum_{k=1}^n \mu(A_k)$.
  \item \textbf{Monotonluk:} $A \subseteq B \Rightarrow \mu(A) \leq \mu(B)$.
  \item \textbf{Alt-toplamsallık:} $\mu\!\left(\bigcup_{n=1}^\infty A_n\right)
        \leq \sum_{n=1}^\infty \mu(A_n)$.
  \item \textbf{Aşağıdan süreklilik:} $A_n \nearrow A \Rightarrow \mu(A_n) \nearrow \mu(A)$.
  \item \textbf{Yukarıdan süreklilik:} $A_n \searrow A$ ve $\mu(A_1) < \infty$
        $\Rightarrow \mu(A_n) \searrow \mu(A)$.
\end{enumerate}
\end{teorem}

\begin{proof}
\textbf{(i)} $A_{n+1} = A_{n+2} = \cdots = \emptyset$ koyup $\sigma$-toplamsallığı uygulayın;
$\mu(\emptyset) = 0$ kullanın.

\textbf{(ii)} $B = A \cup (B \setminus A)$ ayrık birleşim; (i)'den
$\mu(B) = \mu(A) + \mu(B \setminus A) \geq \mu(A)$.

\textbf{(iii)} $B_n = A_n \setminus \bigcup_{k=1}^{n-1} A_k$ tanımlayın.
$(B_n)$ ikili ayrık, $\bigcup B_n = \bigcup A_n$ ve $B_n \subseteq A_n$.
$\sigma$-toplamsallık ve monotonluktan:
$\mu(\bigcup A_n) = \sum \mu(B_n) \leq \sum \mu(A_n)$.

\textbf{(iv)} $B_1 = A_1$, $B_n = A_n \setminus A_{n-1}$ ($n \geq 2$) koyun.
$(B_n)$ ikili ayrık ve $A_n = \bigcup_{k=1}^n B_k$, $A = \bigcup_{k=1}^\infty B_k$.
$\sigma$-toplamsallıktan:
\[
  \mu(A) = \sum_{k=1}^\infty \mu(B_k) = \lim_{n\to\infty} \sum_{k=1}^n \mu(B_k) = \lim_{n\to\infty} \mu(A_n).
\]

\textbf{(v)} $C_n = A_1 \setminus A_n$ uygulayın; $C_n \nearrow A_1 \setminus A$.
(iv)'ten $\mu(C_n) \nearrow \mu(A_1 \setminus A) = \mu(A_1) - \mu(A)$.
$\mu(A_n) = \mu(A_1) - \mu(C_n) \to \mu(A)$.
($\mu(A_1) < \infty$ koşulu farkın tanımlı olması için gereklidir.)
\end{proof}

\begin{hataanalizi}[Yukarıdan süreklilik için sonluluk koşulu]
$\mu(A_1) < \infty$ koşulu kaldırılamaz. Karşı örnek: $\Omega = \mathbb{N}$,
$\mu$ sayma ölçüsü ($\mu(\{k\}) = 1$ her $k$ için), $A_n = \{n, n+1, n+2, \ldots\}$.
O zaman $A_n \searrow \emptyset$ ama $\mu(A_n) = \infty$ her $n$ için; $\mu(\emptyset) = 0$'a yakınsamaz.
\end{hataanalizi}

\begin{ornek}[Sayma ölçüsü ve Dirac ölçüsü]\label{ex:saymadirac}
\begin{enumerate}[label=(\alph*)]
  \item \textbf{Sayma ölçüsü} \emph{(counting measure)}: $\Omega$ herhangi,
        $\mathcal{F} = \mathcal{P}(\Omega)$,
        $\mu(A) = |A|$ (eleman sayısı, sonsuz ise $+\infty$).

  \item \textbf{Dirac ölçüsü} \emph{(Dirac measure)}: $x_0 \in \Omega$ sabit,
        \[
          \delta_{x_0}(A) = \gostf{A}(x_0) =
          \begin{cases} 1 & x_0 \in A, \\ 0 & x_0 \notin A. \end{cases}
        \]
        $\delta_{x_0}$ bir olasılık ölçüsüdür. Tüm ölçü $x_0$ noktasında yoğunlaşmıştır.
\end{enumerate}
\end{ornek}

\begin{mikroozet}[§1.2 özeti]
Ölçü: $\mu(\emptyset)=0$ ve $\sigma$-toplamsallık. Olasılık ölçüsü: ölçü + $\mu(\Omega)=1$.
Temel özellikler: monotonluk, alt-toplamsallık, aşağıdan/yukarıdan süreklilik
(sonluluk koşuluyla). Örnekler: sayma ölçüsü, Dirac ölçüsü, Lebesgue ölçüsü.
\end{mikroozet}


% =====================================================================
\section{Dış Ölçü ve Carathéodory'nin Genişleme Teoremi}
% =====================================================================

\begin{motivasyon}[§1.3]
Lebesgue ölçüsünü nasıl kurarız? ``Uzunluk'' fonksiyonunu aralıklar üzerinde
tanımlamak kolay: $\lambda((a,b)) = b-a$. Ama bu tanımı tüm Borel kümelerine
nasıl genişletiriz? Carathéodory'nin yaklaşımı: önce tüm $\mathcal{P}(\mathbb{R})$
üzerinde bir ``dış ölçü'' tanımla; ardından ölçülebilir kümeleri seç.
\end{motivasyon}

\begin{tanim}[Dış Ölçü]\label{def:dismolcu}
$\mu^*: \mathcal{P}(\Omega) \to [0,+\infty]$ fonksiyonuna \textbf{dış ölçü}
\emph{(outer measure)} denir; eğer:
\begin{enumerate}[label=(\alph*)]
  \item $\mu^*(\emptyset) = 0$,
  \item $A \subseteq B \Rightarrow \mu^*(A) \leq \mu^*(B)$,
  \item $\mu^*\!\left(\bigcup_{n=1}^\infty A_n\right) \leq \sum_{n=1}^\infty \mu^*(A_n)$.
\end{enumerate}
\end{tanim}

\begin{tanim}[Carathéodory Ölçülebilirliği]\label{def:caratheodory}
$E \subseteq \Omega$, $\mu^*$'a göre \textbf{Carathéodory-ölçülebilir} adını alır;
eğer her $A \subseteq \Omega$ için
\[
  \mu^*(A) = \mu^*(A \cap E) + \mu^*(A \cap E^c).
\]
Carathéodory-ölçülebilir kümelerin ailesi $\mathcal{M}^*$ ile gösterilir.
\end{tanim}

\begin{teorem}[Carathéodory'nin Genişleme Teoremi]\label{thm:caratheodory}
$\mu^*$ bir dış ölçü olsun. O zaman:
\begin{enumerate}[label=(\roman*)]
  \item $\mathcal{M}^*$ bir $\sigma$-cebiriyidir.
  \item $\mu^*\!\restriction_{\mathcal{M}^*}$ bir ölçüdür.
\end{enumerate}
\end{teorem}

\begin{proof}
\textit{Fikir:} Carathéodory koşulu, $E$'nin her $A$'yı ``mükemmel biçimde''
iki parçaya böldüğünü söyler; bu özellik $\sigma$-cebiri aksiyomlarıyla uyumludur.

\medskip

\textbf{$\mathcal{M}^*$ cebirdir (sonlu kapatma).}

\textit{Tümleme:} $E \in \mathcal{M}^*$ ise $E^c \in \mathcal{M}^*$: tanım simetrik.

\textit{Sonlu birleşim:} $E, F \in \mathcal{M}^*$, $A \subseteq \Omega$ olsun.
$E$'nin Carathéodory koşulunu $A \cap F^c$'ye uygula:
\[
  \mu^*(A \cap F^c) = \mu^*(A \cap F^c \cap E) + \mu^*(A \cap F^c \cap E^c).
\]
$F$'nin koşulunu $A$'ya uygula: $\mu^*(A) = \mu^*(A \cap F) + \mu^*(A \cap F^c)$.
Toplarsak ve $A \cap (E \cup F)^c = A \cap E^c \cap F^c$ göz önünde tutulursa,
$\mu^*(A) = \mu^*(A \cap (E \cup F)) + \mu^*(A \cap (E \cup F)^c)$.

\medskip

\textbf{$\mathcal{M}^*$ sayılabilir birleşim altında kapalıdır ve $\sigma$-toplamsallık.}

$(E_n) \subseteq \mathcal{M}^*$ ikili ayrık, $S_n = \bigcup_{k=1}^n E_k$, $S = \bigcup_n E_n$.
Tümevarımla $S_n \in \mathcal{M}^*$. Her $A$ için:
\[
  \mu^*(A \cap S_n) = \sum_{k=1}^n \mu^*(A \cap E_k).
\]
(Bu eşitlik $E_k$'nin Carathéodory koşulunun $A \cap S_k$'ya uygulanmasıyla tümevarımla elde edilir.)
$S_n \in \mathcal{M}^*$ koşulunu $A$'ya uygula ve $n \to \infty$ al:
\[
  \mu^*(A) = \mu^*(A \cap S_n) + \mu^*(A \cap S_n^c)
  \geq \sum_{k=1}^n \mu^*(A \cap E_k) + \mu^*(A \cap S^c).
\]
$n \to \infty$: $\mu^*(A) \geq \sum_k \mu^*(A \cap E_k) + \mu^*(A \cap S^c)
\geq \mu^*(A \cap S) + \mu^*(A \cap S^c) \geq \mu^*(A)$.
Tüm eşitsizlikler eşitlik verir; $S \in \mathcal{M}^*$ ve $\sigma$-toplamsallık.
\end{proof}

\begin{mikroozet}[§1.3 özeti]
Dış ölçü tüm kümelerde tanımlı ama yalnızca $\sigma$-alt-toplamsaldır.
Carathéodory kriteri ölçülebilir kümeleri seçer; bu kümelerin oluşturduğu
aile bir $\sigma$-cebiriyidir ve dış ölçünün kısıtlaması tam bir ölçüdür.
\end{mikroozet}


% =====================================================================
\section{Lebesgue Ölçüsü}
% =====================================================================

\begin{tanim}[Lebesgue Dış Ölçüsü]\label{def:lebesgue_dis}
$A \subseteq \mathbb{R}$ için \textbf{Lebesgue dış ölçüsü}:
\[
  \lambda^*(A) \coloneqq \inf\left\{
    \sum_{n=1}^\infty (b_n - a_n) :
    A \subseteq \bigcup_{n=1}^\infty (a_n, b_n)
  \right\}.
\]
\end{tanim}

\begin{teorem}[Lebesgue Ölçüsünün Varlığı]\label{thm:lebesgue_varlik}
$\lambda^*$ bir dış ölçüdür ve $\lambda^*((a,b)) = b - a$ her $a < b$ için.
Carathéodory teoremiyle, $\lambda^*$'nin kısıtlaması
$\mathcal{L} \supseteq \mathcal{B}(\mathbb{R})$ üzerinde bir ölçü verir.
Bu ölçüye \textbf{Lebesgue ölçüsü} $\lambda$ denir.
\end{teorem}

\begin{proof}[İspat (özet)]
Dış ölçü aksiyomlarının doğruluğu $\lambda^*$ tanımından açıktır.
$\lambda^*((a,b)) = b-a$: $\geq$ yönü $(a,b)$'nin kendisiyle örtülmesinden;
$\leq$ yönü Heine-Borel teoremiyle (her açık örtünün sonlu alt örtüsü var)
sağlanır; detay için Billingsley~\cite{billingsley1995} §3 bakınız.
Borel kümelerinin Carathéodory-ölçülebilir olduğu doğrulanır.
\end{proof}

\begin{onerme}[Lebesgue Ölçüsünün Özellikleri]\label{prop:lebesgueozell}
Lebesgue ölçüsü $\lambda$:
\begin{enumerate}[label=(\roman*)]
  \item \textbf{Ötelemede değişmez:} $\lambda(A + x) = \lambda(A)$ her $x \in \mathbb{R}$,
        $A \in \mathcal{L}$ için.
  \item \textbf{Ölçeklemede:} $\lambda(cA) = |c|\,\lambda(A)$ her $c \in \mathbb{R}$.
  \item \textbf{$\sigma$-sonlu:} $\mathbb{R} = \bigcup_n (-n,n)$ ve $\lambda((-n,n)) = 2n < \infty$.
  \item \textbf{Tam:} $\lambda(N) = 0$ ve $A \subseteq N$ ise $A \in \mathcal{L}$ ve $\lambda(A) = 0$.
\end{enumerate}
\end{onerme}

\begin{karsiornek}[Vitali kümesi --- Lebesgue-ölçülemeyen küme]
Seçim aksiyomu kabul edilirse, $[0,1]$'de Lebesgue-ölçülemeyen kümeler
mevcuttur (\textbf{Vitali kümesi}). İnşaat: $[0,1]$'deki sayıları
$\sim$ ile $x \sim y \iff x - y \in \mathbb{Q}$ denk bağıntısıyla sınıflandırın.
Her denk sınıftan bir eleman seçerek $V$ oluşturun. $V$ ölçülebilir değildir:
ölçülebilir olsaydı ya $\lambda(V) = 0$ (ama $[0,1] \subseteq \bigcup_{q \in \mathbb{Q} \cap [-1,1]} (V+q)$,
alt-toplamsallıkla çelişir) ya da $\lambda(V) > 0$ (ama $\lambda([0,2]) \geq \sum_{q} \lambda(V+q)$,
ötelemede değişmezlikle $= \infty$; çelişki).
\end{karsiornek}

\begin{mikroozet}[§1.4 özeti]
Lebesgue ölçüsü aralıkların uzunluğunu tüm Borel (ve daha geniş Lebesgue)
kümelerine genişletir. Öteleme değişmezliği, tamlık ve $\sigma$-sonluluk
temel özellikleridir. Ölçüm teorisinin zorunlu kıldığı Seçim Aksiyomu,
ölçülemeyen kümelerin varlığını garanti eder.
\end{mikroozet}


% =====================================================================
\section{Ölçülebilir Fonksiyonlar}
% =====================================================================

\begin{tanimgerekce}
Rasgele değişken, sezgisel olarak ``rastlantıya bağlı sayısal bir sonuç''tur.
Ama hangi fonksiyonlar rasgele değişken olarak kabul edilebilir? Cevap:
olayları olaylara götürenler, yani ters görüntüsü $\sigma$-cebirini koruyanlar.
\end{tanimgerekce}

\begin{tanim}[Ölçülebilir Fonksiyon]\label{def:olculebilfon}
$(\Omega, \mathcal{F})$ ve $(S, \mathcal{G})$ ölçülebilir uzaylar,
$f: \Omega \to S$ bir fonksiyon olsun.
$f$'ye $\mathcal{F}/\mathcal{G}$-\textbf{ölçülebilir} denir;
eğer her $B \in \mathcal{G}$ için $f^{-1}(B) \in \mathcal{F}$.

Özel olarak $S = \mathbb{R}$, $\mathcal{G} = \mathcal{B}(\mathbb{R})$ ise
$f$ kısaca $\mathcal{F}$-\textbf{ölçülebilir} adını alır.
\end{tanim}

\begin{teorem}[Ölçülebilirlik Karakterizasyonları]\label{thm:olculebichar}
$f: (\Omega,\mathcal{F}) \to (\mathbb{R}, \mathcal{B}(\mathbb{R}))$ için
aşağıdakiler denktir:
\begin{enumerate}[label=(\roman*)]
  \item $f$ ölçülebilirdir.
  \item Her $b \in \mathbb{R}$ için $\{f \leq b\} \coloneqq \{\omega : f(\omega) \leq b\} \in \mathcal{F}$.
  \item Her $b \in \mathbb{R}$ için $\{f < b\} \in \mathcal{F}$.
  \item Her $b \in \mathbb{R}$ için $\{f \geq b\} \in \mathcal{F}$.
  \item Her $b \in \mathbb{R}$ için $\{f > b\} \in \mathcal{F}$.
\end{enumerate}
\end{teorem}

\begin{proof}
(i) $\Rightarrow$ (ii): $(-\infty,b] \in \mathcal{B}(\mathbb{R})$
(Teorem~\ref{thm:boreldenk}(v)), dolayısıyla $f^{-1}((-\infty,b]) \in \mathcal{F}$.

(ii) $\Rightarrow$ (i): $\mathcal{B}(\mathbb{R}) = \sigma(\{(-\infty,b] : b \in \mathbb{R}\})$.
Koşul (ii), $\{(-\infty,b]\}$ ailesinin ters görüntülerinin $\mathcal{F}$'de
olduğunu söyler. $\{B : f^{-1}(B) \in \mathcal{F}\}$ ailesinin bir $\sigma$-cebiri
olduğu görülür (Bölüm~0, Teorem~\ref{thm:tergoriuntu}); bu $\sigma$-cebiri
$\{(-\infty,b]\}$ ailesini içerdiğinden $\mathcal{B}(\mathbb{R})$'yi içerir.

(ii) $\Leftrightarrow$ (iii)---(v) benzer dönüşümlerle: $\{f < b\} = \bigcup_n \{f \leq b - 1/n\}$, vb.
\end{proof}

\begin{onerme}[Ölçülebilir Fonksiyonların Kararlılığı]\label{prop:olculebikarar}
\begin{enumerate}[label=(\roman*)]
  \item $f, g$ ölçülebilir, $c \in \mathbb{R}$ sabit $\Rightarrow$
        $f+g$, $fg$, $cf$, $|f|$, $f \vee g$, $f \wedge g$ ölçülebilir.
  \item $f_n$ ölçülebilir ($n \geq 1$) $\Rightarrow$
        $\sup_n f_n$, $\inf_n f_n$, $\limsup_n f_n$, $\liminf_n f_n$ ölçülebilir.
  \item $f_n \to f$ noktasal ise ve her $f_n$ ölçülebilir ise $f$ ölçülebilir.
\end{enumerate}
\end{onerme}

\begin{proof}
(i) $\{f + g > b\} = \bigcup_{q \in \mathbb{Q}} (\{f > q\} \cap \{g > b-q\})$;
rasyoneller sayılabilir, her küme $\mathcal{F}$'de. Geri kalanlar benzer.

(ii) $\{\sup_n f_n > b\} = \bigcup_n \{f_n > b\} \in \mathcal{F}$.
$\limsup_n f_n = \inf_n \sup_{k \geq n} f_k$; (ii)'nin tekrarlı uygulaması.

(iii) Noktasal limit $\limsup = \liminf$ olduğu durumdur; (ii)'den.
\end{proof}

\begin{tanim}[Basit Fonksiyon]\label{def:basit}
$f: \Omega \to \mathbb{R}$ \textbf{basit} \emph{(simple function)} adını alır; eğer
yalnızca sonlu sayıda değer alıyorsa. Standart gösterim:
\[
  f = \sum_{k=1}^n c_k \gostf{A_k},
  \quad c_k \in \mathbb{R},\quad A_k = f^{-1}(\{c_k\}) \in \mathcal{F},
  \quad A_k \text{ ikili ayrık.}
\]
\end{tanim}

\begin{teorem}[Basit Fonksiyon Yaklaşımı]\label{thm:basityaklasim}
$f: (\Omega, \mathcal{F}) \to [0,\infty]$ ölçülebilir ise, her $n \geq 1$ için
$0 \leq f_n \leq f_{n+1} \leq f$ ve $f_n \nearrow f$ noktasal olan ölçülebilir
basit fonksiyonlar $(f_n)$ mevcuttur.
\end{teorem}

\begin{proof}
$f_n(\omega) = \sum_{k=0}^{n 2^n - 1} \frac{k}{2^n} \gostf{\{k/2^n \leq f < (k+1)/2^n\}}(\omega) + n \gostf{\{f \geq n\}}(\omega)$
tanımlayın. $f_n$ ölçülebilir (her $A_k \in \mathcal{F}$), basit (sonlu değer alır),
artan ($f_n \leq f_{n+1}$), ve $f_n(\omega) \nearrow f(\omega)$ her $\omega$ için
($f(\omega) < \infty$ ise yakınsama açık; $f(\omega) = \infty$ ise $f_n(\omega) = n \to \infty$).
\end{proof}

\begin{mikroozet}[§1.5 özeti]
Ölçülebilir fonksiyon: ters görüntüsü $\sigma$-cebirini koruyan fonksiyon.
Karakterizasyon: $\{f \leq b\} \in \mathcal{F}$ her $b$ için yeterli ve gerekli.
Basit fonksiyonlar ölçülebilir fonksiyonları aşağıdan yaklaşır: integral
teorisinin temel aracı.
\end{mikroozet}


% =====================================================================
\section{Lebesgue İntegrali}
% =====================================================================

\begin{motivasyon}[§1.6]
Lebesgue integrali üç aşamada tanımlanır: önce basit fonksiyonlar için
(sezgisel), sonra negatif olmayan fonksiyonlar için (supremum yoluyla),
son olarak işaretli fonksiyonlar için ($f = f^+ - f^-$). Her aşama
bir öncekine dayanır.
\end{motivasyon}

\begin{tanim}[Lebesgue İntegrali --- Üç Aşama]\label{def:lebesgueintegral}
$(\Omega, \mathcal{F}, \mu)$ ölçüm uzayı.

\medskip

\textbf{Aşama 1.} $f = \sum_{k=1}^n c_k \gostf{A_k}$ ($c_k \geq 0$) basit:
\[
  \int f \, d\mu \coloneqq \sum_{k=1}^n c_k \, \mu(A_k) \in [0,\infty].
\]

\textbf{Aşama 2.} $f: \Omega \to [0,\infty]$ ölçülebilir:
\[
  \int f \, d\mu \coloneqq \sup\left\{ \int g \, d\mu :
    g \text{ basit, } 0 \leq g \leq f \right\}.
\]

\textbf{Aşama 3.} $f: \Omega \to \mathbb{R}$ ölçülebilir. $f^+ = \max(f,0)$,
$f^- = \max(-f,0)$. $f$ \textbf{$\mu$-integrallenebilir} adını alır; eğer
$\int f^+ \, d\mu < \infty$ ve $\int f^- \, d\mu < \infty$. O zaman:
\[
  \int f \, d\mu \coloneqq \int f^+ \, d\mu - \int f^- \, d\mu.
\]
$L^1(\mu) \coloneqq \{f \text{ ölçülebilir} : \int |f| \, d\mu < \infty\}$.
\end{tanim}

\begin{notkutu}
$A$ üzerinde integral: $\int_A f \, d\mu \coloneqq \int f \gostf{A} \, d\mu$.
Olasılık teorisinde $\mu = \mathbb{P}$: $\int f \, d\mathbb{P} = \mathbb{E}[f]$.
\end{notkutu}

\begin{teorem}[Lebesgue İntegralin Özellikleri]\label{thm:lebesgueozell}
$f, g \in L^1(\mu)$, $c \in \mathbb{R}$:
\begin{enumerate}[label=(\roman*)]
  \item \textbf{Doğrusallık:} $\int (cf + g) \, d\mu = c \int f \, d\mu + \int g \, d\mu$.
  \item \textbf{Monotonluk:} $f \leq g$ n.k. $\Rightarrow \int f \, d\mu \leq \int g \, d\mu$.
  \item \textbf{Modül eşitsizliği:} $\left|\int f \, d\mu\right| \leq \int |f| \, d\mu$.
  \item \textbf{Sıfır ölçülü küme:} $\mu(A) = 0 \Rightarrow \int_A f \, d\mu = 0$.
  \item \textbf{N.k. eşitlik:} $f = g$ n.k. $\Rightarrow \int f \, d\mu = \int g \, d\mu$.
\end{enumerate}
\end{teorem}

\begin{proof}
(i) Basit fonksiyonlar için doğrusallık tanımdan; genel durum Teorem~\ref{thm:basityaklasim}
ve MYT (§1.7) ile taşınır.
(ii) $g - f \geq 0$ n.k.; Aşama 2 tanımdan $\int(g-f) \geq 0$.
(iii) $|f| = f^+ + f^-$, üçgen eşitsizliği.
(iv), (v) Aşama 1 tanımdan; $\mu(A) = 0$ ise $c_k \mu(A_k \cap A) = 0$.
\end{proof}

\begin{mikroozet}[§1.6 özeti]
Lebesgue integrali: basit fonksiyonlar $\to$ negatif olmayan $\to$ işaretli.
$L^1(\mu)$: integrallenebilir fonksiyonlar. Doğrusallık, monotonluk, n.k.
eşitlik temel özellikleri.
\end{mikroozet}


% =====================================================================
\section{Monoton Yakınsama, Fatou ve Dominante Yakınsama}
% =====================================================================

\begin{motivasyon}[§1.7]
Ne zaman $\lim_n \int f_n \, d\mu = \int \lim_n f_n \, d\mu$ yazabiliriz?
Riemann teorisinde bu çok kısıtlıdır. Lebesgue teorisinin gücü burada ortaya
çıkar: üç büyük geçiş teoremi bu soruyu kapsamlı biçimde yanıtlar.
\end{motivasyon}

\begin{teorem}[Monoton Yakınsama Teoremi (MYT)]\label{thm:myt}
$f_n: (\Omega,\mathcal{F},\mu) \to [0,\infty]$ ölçülebilir ve $f_n \nearrow f$
noktasal ise:
\[
  \int f \, d\mu = \lim_{n\to\infty} \int f_n \, d\mu.
\]
\end{teorem}

\begin{proof}
\textit{Fikir:} $\int f_n$ artan ve $\leq \int f$; yeterliyse limit $\int f$'ye eşit.

\medskip

\textbf{Adım 1.} $f_n \leq f$ ve monotonluktan $\int f_n \leq \int f$; dolayısıyla
$L \coloneqq \lim_n \int f_n \leq \int f$.

\textbf{Adım 2.} $L \geq \int f$ göstereceğiz. $0 < c < 1$ ve $g \leq f$ basit
fonksiyon olsun. $E_n = \{f_n \geq c \cdot g\}$ tanımla; $f_n \nearrow f \geq g$'den
$E_n \nearrow \Omega$.

$\int f_n \, d\mu \geq \int_{E_n} f_n \, d\mu \geq c \int_{E_n} g \, d\mu$.

$g = \sum_k a_k \gostf{A_k}$ basit; aşağıdan süreklilikten (Teorem~\ref{thm:olcuozellik}(iv)):
$\int_{E_n} g \, d\mu = \sum_k a_k \mu(A_k \cap E_n) \to \sum_k a_k \mu(A_k) = \int g \, d\mu$.

Dolayısıyla $L \geq c \int g \, d\mu$. $c \to 1^-$ ve ardından
$g \nearrow f$ (Teorem~\ref{thm:basityaklasim}) ile $L \geq \int f$.
\end{proof}

\begin{sonuc}[MYT --- Seriler]\label{cor:mytseriler}
$f_n \geq 0$ ölçülebilir ise $\int \sum_{n=1}^\infty f_n \, d\mu = \sum_{n=1}^\infty \int f_n \, d\mu$.
\end{sonuc}

\begin{proof}
$g_N = \sum_{n=1}^N f_n \nearrow \sum_n f_n$; MYT uygula.
\end{proof}

\begin{teorem}[Fatou Lemması]\label{thm:fatou}
$f_n \geq 0$ ölçülebilir ise:
\[
  \int \liminf_{n\to\infty} f_n \, d\mu \leq \liminf_{n\to\infty} \int f_n \, d\mu.
\]
\end{teorem}

\begin{proof}
$g_n = \inf_{k \geq n} f_k$ tanımla; $g_n \leq f_n$ ve $g_n \nearrow \liminf_n f_n$.
MYT: $\int \liminf_n f_n = \lim_n \int g_n$.
$g_n \leq f_k$ her $k \geq n$ için; dolayısıyla
$\int g_n \leq \inf_{k \geq n} \int f_k$.
$\lim_n \int g_n \leq \lim_n \inf_{k \geq n} \int f_k = \liminf_n \int f_n$.
\end{proof}

\begin{hataanalizi}[Fatou'da eşitsizlik keskin olabilir]
Eşitlik her zaman sağlanmaz. $f_n = \gostf{[n,n+1]}$ alalım; $f_n \to 0$ noktasal,
$\int f_n = 1$ her $n$ için. $\int \liminf f_n = 0 < 1 = \liminf \int f_n$.
\end{hataanalizi}

\begin{teorem}[Dominante Yakınsama Teoremi (DYT)]\label{thm:dyt}
$f_n \to f$ n.k. ve $|f_n| \leq g$ n.k. (her $n$ için), $g \in L^1(\mu)$ ise:
\[
  \lim_{n\to\infty} \int f_n \, d\mu = \int f \, d\mu
  \quad \text{ve} \quad
  \lim_{n\to\infty} \int |f_n - f| \, d\mu = 0.
\]
\end{teorem}

\begin{proof}
\textit{Fikir:} $g + f_n \geq 0$ ve $g - f_n \geq 0$; Fatou uygulanabilir.

\medskip

$g + f_n \geq 0$; Fatou:
$\int (g + f) \leq \liminf_n \int (g + f_n) = \int g + \liminf_n \int f_n$.
$\int g < \infty$'dan $\int f \leq \liminf_n \int f_n$.

$g - f_n \geq 0$; Fatou:
$\int (g - f) \leq \liminf_n \int (g - f_n) = \int g - \limsup_n \int f_n$.
$\limsup_n \int f_n \leq \int f$.

İkisi birden: $\limsup_n \int f_n \leq \int f \leq \liminf_n \int f_n$,
yani $\lim_n \int f_n = \int f$.

$L^1$ yakınsaması: $|f_n - f| \leq 2g \in L^1$; noktasal $\to 0$; DYT yeniden uygulanır.
\end{proof}

\begin{ornek}[DYT uygulaması]\label{ex:dytuyg}
$f_n(x) = \dfrac{nx}{1 + n^2 x^2}$, $x \in [0,1]$, Lebesgue ölçüsüyle.
$f_n(x) \to 0$ her $x > 0$ için (ve $f_n(0) = 0$).
$|f_n(x)| \leq \frac{1}{2}$ (AM-GM: $nx \cdot 1 \leq (nx + 1/(nx))/2$... daha basit:
$f_n'(x) = 0$ çözümü $x = 1/n$; $f_n(1/n) = 1/2$). Dolayısıyla $g(x) = 1/2 \in L^1([0,1])$.
DYT: $\int_0^1 f_n \, dx \to \int_0^1 0 \, dx = 0$.
\end{ornek}

\begin{mikroozet}[§1.7 özeti]
MYT: artan + negatif olmayan $\Rightarrow$ limit içine girebilir.
Fatou: negatif olmayan $\Rightarrow$ $\int \liminf \leq \liminf \int$.
DYT: integrallenebilir dom. fonksiyon varsa noktasal yakınsama $\Rightarrow$ $L^1$ yakınsaması.
\end{mikroozet}


% =====================================================================
\section{Ürün Ölçüleri ve Fubini-Tonelli Teoremi}
% =====================================================================

\begin{motivasyon}[§1.8]
İki rasgele değişkenin ortak dağılımını hesaplamak, $\mathbb{R}^2$ üzerinde
bir integral almak gerektirir. Fubini teoremi bu iki boyutlu integrali
tek boyutlu integrallerin ardışık uygulamasına indirgeyerek hesaplamayı
mümkün kılar.
\end{motivasyon}

\begin{tanim}[Ürün $\sigma$-Cebiri]\label{def:urun_sigma}
$(\Omega_1, \mathcal{F}_1)$ ve $(\Omega_2, \mathcal{F}_2)$ ölçülebilir uzaylar.
\[
  \mathcal{F}_1 \otimes \mathcal{F}_2 \coloneqq
  \sigma\bigl(\{A_1 \times A_2 : A_1 \in \mathcal{F}_1, A_2 \in \mathcal{F}_2\}\bigr)
\]
$\Omega_1 \times \Omega_2$ üzerinde \textbf{ürün $\sigma$-cebiriyidir}.
$A_1 \times A_2$ biçimindeki kümelere \textbf{dikdörtgen (measurable rectangle)} denir.
\end{tanim}

\begin{onerme}
$\mathcal{B}(\mathbb{R}) \otimes \mathcal{B}(\mathbb{R}) = \mathcal{B}(\mathbb{R}^2)$.
\end{onerme}

\begin{proof}
$\mathbb{R}^2$'deki açık dikdörtgenler $(a_1,b_1) \times (a_2,b_2)$ hem
$\mathcal{B}(\mathbb{R}) \otimes \mathcal{B}(\mathbb{R})$'de hem de
$\mathcal{B}(\mathbb{R}^2)$'de; karşılıklı içerme verir.
\end{proof}

\begin{teorem}[Ürün Ölçüsünün Varlığı ve Tekliği]\label{thm:urunoelcu}
$(\Omega_1, \mathcal{F}_1, \mu_1)$ ve $(\Omega_2, \mathcal{F}_2, \mu_2)$ $\sigma$-sonlu
ölçüm uzayları olsun. $\Omega_1 \times \Omega_2$ üzerinde tek bir $\mu_1 \otimes \mu_2$
ölçüsü mevcuttur; öyle ki her dikdörtgen için:
\[
  (\mu_1 \otimes \mu_2)(A_1 \times A_2) = \mu_1(A_1) \cdot \mu_2(A_2).
\]
\end{teorem}

\begin{proof}[İspat (özet)]
Dikdörtgenler bir $\pi$-sistemi oluşturur. Carathéodory genişlemesi
$\mu_1 \otimes \mu_2$'yi inşa eder. $\sigma$-sonluluk tekliği garantilemek
için gereklidir (Dynkin'in $\pi$-$\lambda$ argümanı).
Ayrıntı için Klenke~\cite{klenke2014} §14 bakınız.
\end{proof}

\begin{teorem}[Tonelli Teoremi]\label{thm:tonelli}
$f: (\Omega_1 \times \Omega_2, \mathcal{F}_1 \otimes \mathcal{F}_2) \to [0,\infty]$
ölçülebilir olsun. O zaman:
\begin{enumerate}[label=(\roman*)]
  \item $\omega_1 \mapsto \int f(\omega_1, \omega_2) \, d\mu_2(\omega_2)$ ve
        $\omega_2 \mapsto \int f(\omega_1, \omega_2) \, d\mu_1(\omega_1)$
        ölçülebilir fonksiyonlardır.
  \item
        \[
          \int f \, d(\mu_1 \otimes \mu_2)
          = \int\!\!\int f(\omega_1,\omega_2) \, d\mu_2(\omega_2) \, d\mu_1(\omega_1)
          = \int\!\!\int f(\omega_1,\omega_2) \, d\mu_1(\omega_1) \, d\mu_2(\omega_2).
        \]
\end{enumerate}
\end{teorem}

\begin{teorem}[Fubini Teoremi]\label{thm:fubini}
$f \in L^1(\mu_1 \otimes \mu_2)$ olsun. O zaman:
\begin{enumerate}[label=(\roman*)]
  \item $\mu_1$-n.k. her $\omega_1$ için $\omega_2 \mapsto f(\omega_1, \omega_2) \in L^1(\mu_2)$.
  \item Tonelli teoremindeki tekrarlı integrasyon formülü geçerlidir.
\end{enumerate}
\end{teorem}

\begin{proof}[İspat (özet)]
Tonelli teoremini $f^+$ ve $f^-$'ya ayrı ayrı uygulayın.
$f \in L^1$ olduğundan $\int f^+ < \infty$ ve $\int f^- < \infty$;
$\omega_1$-neredeyse her yerde tekil integraller de sonludur.
\end{proof}

\begin{hataanalizi}[$L^1$ koşulu olmadan Fubini başarısız olabilir]
$f(x,y) = \frac{x^2 - y^2}{(x^2+y^2)^2}$ üzerinde $[0,1]^2$ Lebesgue ölçüsü.
$\int_0^1 \int_0^1 f \, dx \, dy = \pi/4$ ama $\int_0^1 \int_0^1 f \, dy \, dx = -\pi/4$.
$f \notin L^1([0,1]^2)$; Fubini koşulu sağlanmaz.
\end{hataanalizi}

\begin{ornek}[Fubini --- somut uygulama]
$(\mathbb{R}^2, \mathcal{B}(\mathbb{R}^2), \lambda^2)$ üzerinde
$f(x,y) = e^{-x^2 - y^2}$ integralini hesaplayın.
\[
  \int_{\mathbb{R}^2} e^{-x^2-y^2} \, d\lambda^2
  = \int_{-\infty}^\infty e^{-x^2} \, dx \cdot \int_{-\infty}^\infty e^{-y^2} \, dy
  = \left(\int_{-\infty}^\infty e^{-x^2} \, dx\right)^{\!2}.
\]
Kutupsal koordinatlarla $\int_{\mathbb{R}^2} e^{-r^2} r \, dr \, d\theta = \pi$;
dolayısıyla $\int_{-\infty}^\infty e^{-x^2} \, dx = \sqrt{\pi}$.
Bu, Gaussian integralinin en zarif hesabıdır.
\end{ornek}

\begin{mikroozet}[§1.8 özeti]
Ürün $\sigma$-cebiri: dikdörtgenler tarafından üretilen. Ürün ölçüsü dikdörtgenlerde
çarpımsal. Tonelli: negatif olmayan için sıra serbesttir. Fubini: $L^1$'de sıra serbesttir.
$L^1$ koşulu olmadan sıra değiştirme yanlış sonuç verebilir.
\end{mikroozet}


% =====================================================================
\section{$L^p$ Uzayları ve İntegral Eşitsizlikleri}
\label{sec:lp_uzaylari}
% =====================================================================

\begin{tanim}[$L^p$ Uzayı]
\label{def:lp_uzay}
$(\Omega, \mathcal{F}, \mu)$ ölçüm uzayı, $1 \leq p < \infty$.
\[
  L^p(\mu) = \{f \text{ ölçülebilir} : \|f\|_p := \left(\int |f|^p \, d\mu\right)^{1/p} < \infty\}.
\]
$L^\infty(\mu)$: $\|f\|_\infty = \inf\{M : \mu(\{|f| > M\}) = 0\}$ (esansiyel üst sınır).
Teknik olarak $L^p$ sınıfları: eşdeğer fonksiyonlar (n.k.\ eşit) aynı sınıfa.
\end{tanim}

\begin{teorem}[Hölder Eşitsizliği]
\label{thm:holder}
$1/p + 1/q = 1$ ($p,q > 1$), $f \in L^p$, $g \in L^q$. O zaman $fg \in L^1$ ve
\[
  \int |fg| \, d\mu \leq \|f\|_p \, \|g\|_q.
\]
Özel durum $p = q = 2$: Cauchy-Schwarz eşitsizliği $\int|fg| \leq \|f\|_2\|g\|_2$.
\end{teorem}

\begin{proof}
Young eşitsizliği: $ab \leq a^p/p + b^q/q$ ($a,b \geq 0$, $1/p+1/q=1$).
$a = |f|/\|f\|_p$, $b = |g|/\|g\|_q$ koy; $\mu$-integral al:
$\int |fg|/(\|f\|_p\|g\|_q) \leq \int |f|^p/(p\|f\|_p^p) + \int|g|^q/(q\|g\|_q^q) = 1/p + 1/q = 1$.
\end{proof}

\begin{teorem}[Minkowski Eşitsizliği]
\label{thm:minkowski}
$1 \leq p \leq \infty$, $f, g \in L^p$. O zaman $f + g \in L^p$ ve
\[
  \|f + g\|_p \leq \|f\|_p + \|g\|_p.
\]
\end{teorem}

\begin{proof}
$p=1$: üçgen eşitsizliği, $|f+g| \leq |f| + |g|$.
$1 < p < \infty$: $\|f+g\|_p^p = \int|f+g|^p \leq \int|f+g|^{p-1}|f| + \int|f+g|^{p-1}|g|$.
Her terime Hölder ($|f+g|^{p-1} \in L^q$ çünkü $(p-1)q = p$):
$\leq \|f+g\|_p^{p-1}(\|f\|_p + \|g\|_p)$. $\|f+g\|_p^{p-1}$'e bölme ile sonuç.
\end{proof}

\begin{teorem}[$L^p$ Uzayı Tamlığı — Riesz-Fischer]
\label{thm:riesz_fischer}
$L^p(\mu)$ ($1 \leq p \leq \infty$) Minkowski normuyla tam metrize edilebilir uzaydır
(Banach uzayı). $p = 2$: Hilbert uzayı (iç çarpım $\langle f,g\rangle = \int fg$).
\end{teorem}

\begin{teorem}[$L^p$ Dahil Olma İlişkisi]
\label{thm:lp_dahil}
$\mu(\Omega) < \infty$ ise $1 \leq p \leq q \leq \infty$ için $L^q \subseteq L^p$.
\end{teorem}

\begin{proof}
Hölder: $\|f\|_p^p = \int|f|^p \cdot 1 \leq \|f^p\|_{q/p} \cdot \|1\|_{q/(q-p)} = \|f\|_q^p \cdot \mu(\Omega)^{1-p/q}$.
\end{proof}

\begin{ornek}[$L^p$ Farklılıkları]
$\Omega = (0,1)$, $\lambda$ Lebesgue. $f(x) = x^{-\alpha}$ ($0 < \alpha < 1$):
$\int_0^1 x^{-\alpha p} \, dx < \infty \Leftrightarrow \alpha p < 1 \Leftrightarrow p < 1/\alpha$.
Yani $f \in L^p$ ancak $p < 1/\alpha$ için. Küçük $p$'ler için $f \in L^p$, büyük $p$'ler için değil.
\end{ornek}

\begin{mikroozet}
$L^p$: $p$-inci mutlak moment sonlu fonksiyonlar. Hölder: $\int|fg| \leq \|f\|_p\|g\|_q$.
Minkowski: $\|f+g\|_p \leq \|f\|_p + \|g\|_p$ (üçgen). $L^p$ Banach; $L^2$ Hilbert.
\end{mikroozet}

% =====================================================================
\section{Radon-Nikodym Teoremi}
\label{sec:radon_nikodym}
% =====================================================================

\begin{tanim}[Mutlak Süreklilik]
\label{def:mutlak_sureklilik}
$\mu$ ve $\nu$ aynı $(\Omega, \mathcal{F})$ üzerinde ölçüler.
$\nu$, $\mu$'ya göre \textbf{mutlak sürekli} ($\nu \ll \mu$) denir, eğer
her $A \in \mathcal{F}$ için $\mu(A) = 0 \Rightarrow \nu(A) = 0$.
\end{tanim}

\begin{ornek}[Mutlak Süreklilik]
$f \geq 0$, $\int f \, d\mu = 1$. $\nu(A) = \int_A f \, d\mu$: o zaman $\nu \ll \mu$.
$\mu(A) = 0 \Rightarrow \int_A f \, d\mu = 0 \Rightarrow \nu(A) = 0$.
\end{ornek}

\begin{teorem}[Radon-Nikodym Teoremi]
\label{thm:radon_nikodym}
$(\Omega, \mathcal{F})$ üzerinde $\sigma$-sonlu $\mu$ ve $\nu$ ölçüleri, $\nu \ll \mu$.
O zaman tek türlü ($\mu$-n.k.) ölçülebilir $f \geq 0$ ($\mu$-n.k.) fonksiyon vardır:
\[
  \nu(A) = \int_A f \, d\mu \quad \text{her } A \in \mathcal{F} \text{ için.}
\]
$f = d\nu/d\mu$: \textbf{Radon-Nikodym türevi} (yoğunluk / Nikodym türevi).
\end{teorem}

\begin{proof}[İspat Taslağı (Hilbert Uzayı Argümanı)]
$\rho = \mu + \nu$. $\ell(g) = \int g \, d\nu$ ($g \in L^2(\rho)$) sınırlı lineer fonksiyonel.
Riesz temsil: $\ell(g) = \int gh \, d\rho$ olan $h \in L^2(\rho)$ var.
$0 \leq h \leq 1$ $\rho$-n.k.\ gösterilebilir. $f = h/(1-h)$ ($h < 1$ üzerinde) $\mu$-n.k.\ tanımlı;
$\nu(A) = \int_A f \, d\mu$ doğrulanır. $f$ benzersizliği: iki RN türevi $f_1, f_2$
ise $\int_A(f_1-f_2)d\mu = 0$ her $A$'da; $f_1 = f_2$ $\mu$-n.k.
\end{proof}

\begin{onerme}[Zincir Kuralı]
$\nu \ll \mu \ll \lambda$ ise $d\nu/d\lambda = (d\nu/d\mu)(d\mu/d\lambda)$ ($\lambda$-n.k.).
\end{onerme}

\begin{ornek}[Olasılıkta Önemi]
$X$ rasgele değişken, $P_X$ görüntü ölçüsü (dağılımı). $P_X \ll \lambda$ ise
RN türevi $f = dP_X/d\lambda$: bu tam olarak $X$'in PDF'sidir.
\emph{Koşullu beklenti} $E[X \mid \mathcal{G}]$ Radon-Nikodym teoreminden doğar
(Bölüm~\ref{chp:beklenti}'de ayrıntılı işlendi).
\end{ornek}

\begin{disiplinlerarasibaglan}
\textbf{İstatistik:} Olabilirlik oranı (likelihood ratio) $d\nu/d\mu$ Radon-Nikodym türevidir;
Neyman-Pearson lemması ve yeterli istatistik teorisi bu çerçevede kurulur.\\
\textbf{Bayesçi çıkarım:} Sonsal dağılım, önsel üzerine olabilirlik oranı ile Radon-Nikodym güncellemesidir.\\
\textbf{Stokastik kalkülüs:} Girsanov teoremi, olasılık ölçüsü değiştirmeyi Radon-Nikodym teoremiyle gerçekleştirir.
\end{disiplinlerarasibaglan}

% =====================================================================
\section{Lebesgue-Stieltjes İntegrali ve Beklenti}
\label{sec:ls_integral}
% =====================================================================

\begin{tanim}[Lebesgue-Stieltjes ölçüsü]\label{def:ls_olcusu}
$F:\mathbb{R}\to[0,1]$ KDF (sağdan sürekli, artan, $F(-\infty)=0$, $F(+\infty)=1$). Carathéodory genişlemesiyle $\mu_F$ tanımlanır: $\mu_F((a,b])=F(b)-F(a)$. Bu \emph{Lebesgue-Stieltjes ölçüsüdür}.
\end{tanim}

\begin{teorem}[Beklenti olarak LS integrali]\label{thm:beklenti_ls}
$X$ rasgele değişken, $F_X$ KDF. Her borelci $g:\mathbb{R}\to\mathbb{R}$ için:
\[
  \Beklenti[g(X)] = \int_\Omega g(X(\omega))\,dP(\omega) = \int_\mathbb{R} g(x)\,d\mu_{F_X}(x)
  = \int_\mathbb{R} g(x)\,dF_X(x).
\]
Sürekli $F_X$'te son integral Riemann-Stieltjes'e indirgenir.
\end{teorem}

\begin{ornek}[Karma dağılım beklentisi]\label{ex:karma_beklenti}
$X$: $p=1/3$ olasılıkla $X=0$ (kesikli kısım), $(1-p)=2/3$ olasılıkla $\Uniform(0,1)$ (sürekli kısım).
$F_X(x)=\frac13\mathbf{1}[x\geq0]+\frac23 x\mathbf{1}[0\leq x<1]+\frac23\mathbf{1}[x\geq1]$.
$\Beklenti[X]=\frac13\cdot0+\frac23\int_0^1x\,dx=\frac23\cdot\frac12=\frac13$.
LS integrali karma durumlarda tek çatı altında birleştirir.
\end{ornek}

\begin{teorem}[Riemann $\leftrightarrow$ Lebesgue]\label{thm:riemann_lebesgue}
$f:[a,b]\to\mathbb{R}$ sınırlı. $f$ Riemann integrallenebilir $\iff$ $f$'nin süreksizlik noktaları $\lambda$-ölçüsü sıfır kümedir. Bu durumda Riemann ve Lebesgue integralleri eşittir.
\end{teorem}

\begin{hataanalizi}
$f(x)=\mathbf{1}_\mathbb{Q}(x)$ Dirichlet fonksiyonu: Riemann integrallenemez (her alt aralıkta rasyonel + irrasyonel var), ama Lebesgue integrali sıfırdır ($\lambda(\mathbb{Q})=0$). Bu, Lebesgue integralinin Riemann'ı kesinlikle kapsadığını gösterir.
\end{hataanalizi}

% =====================================================================
\section{Olasılık Ölçülerinin Zayıf Topolojisi}
\label{sec:zayif_topoloji}
% =====================================================================

\begin{tanim}[Olasılık ölçülerinin zayıf yakınsaması]\label{def:zayif_yakin}
$\mathcal{P}(\mathbb{R})$ olasılık ölçüleri uzayı. $\mu_n\overset{w}{\to}\mu$: her sınırlı, sürekli $f$ için $\int f\,d\mu_n\to\int f\,d\mu$. Bu \emph{zayıf yakınsama} (dağılımda yakınsama ile eşdeğer).
\end{tanim}

\begin{teorem}[Portmanteau teoremi]\label{thm:portmanteau}
$\mu_n\overset{w}{\to}\mu$ aşağıdaki koşulların herhangi biriyle eşdeğerdir:
\begin{enumerate}[(i)]
  \item Her kapalı $F$ için $\limsup\mu_n(F)\leq\mu(F)$.
  \item Her açık $G$ için $\liminf\mu_n(G)\geq\mu(G)$.
  \item Her $\mu$-sınır kümesi $A$ ($\mu(\partial A)=0$) için $\mu_n(A)\to\mu(A)$.
  \item Her sınırlı, Lipschitz $f$ için $\int f\,d\mu_n\to\int f\,d\mu$.
\end{enumerate}
\end{teorem}

\begin{ornek}[Portmanteau uygulaması]\label{ex:portmanteau_uygulama}
$X_n\dto X\sim\Normal(0,1)$. $A=(-\infty,1.96]$: $\mu(\partial A)=\mu(\{1.96\})=0$ ($\Normal$ sürekli). Portmanteau (iii): $P(X_n\leq1.96)\to\Phi(1.96)=0.975$. Normal tablolar bu güvenceyle kullanılabilir.
\end{ornek}

\begin{disiplinlerarasibaglan}
\textbf{Bayes İstatistiği:} Zayıf topoloji, büyük örneklem Bayes posteriorlarının (Bernstein-von Mises) incelenmesinde kullanılır — posteriorun frekansçı normaline zayıf yakınsaması.

\textbf{Optik Ağlar/Kuyruk Teorisi:} Servis süreleri dağılımları zayıf yakınsama çerçevesinde ele alınır; Portmanteau teoremi ağ kapasitesi sınırlarını garanti eder.
\end{disiplinlerarasibaglan}

\begin{mikroozet}
$\nu \ll \mu$: sıfır $\mu$-ölçülü kümeler sıfır $\nu$-ölçülü.
Radon-Nikodym: $\nu \ll \mu$ ve $\sigma$-sonlu $\Rightarrow$ $\nu(A) = \int_A f\,d\mu$ olan $f = d\nu/d\mu$ var.
$f$: koşullu beklenti ve PDF'nin ölçüm teorisi temeli.
LS integrali: karma dağılımları tek çatıda birleştirir; Riemann $\subset$ Lebesgue.
Portmanteau: zayıf yakınsama için 4 eşdeğer koşul; normal tablo kullanımının garantisi.
\end{mikroozet}

% =====================================================================
%  BÖLÜM SONU ÖZETİ
% =====================================================================
\section*{Bölüm Özeti}
\addcontentsline{toc}{section}{Bölüm Özeti}

\begin{bolumozeti}[Bölüm 1 --- Temel Sonuçlar]
\begin{itemize}
  \item \textbf{$\sigma$-cebiri:} (A)~$\Omega \in \mathcal{F}$;
        (B)~tümleme kapalı; (C)~sayılabilir birleşim kapalı.
        $\sigma(\mathcal{E})$: $\mathcal{E}$'yi içeren en küçük.
        $\mathcal{B}(\mathbb{R}) = \sigma(\text{açık aralıklar})$.
  \item \textbf{Ölçü:} $\mu(\emptyset)=0$ ve $\sigma$-toplamsallık.
        Monotonluk, alt-toplamsallık, aşağıdan/yukarıdan süreklilik
        (yukarıdan için $\mu(A_1)<\infty$).
  \item \textbf{Carathéodory:} Dış ölçü $\mu^*$ verilince, Carathéodory-ölçülebilir
        kümeler bir $\sigma$-cebiri oluşturur ve $\mu^*$ bu ailede ölçüdür.
  \item \textbf{Lebesgue:} $\lambda^*$ dış ölçüden $\lambda$ elde edilir;
        $\lambda((a,b))=b-a$, öteleme değişmez, tam ve $\sigma$-sonlu.
  \item \textbf{Ölçülebilir $f$:} $f^{-1}(B) \in \mathcal{F}$ her $B \in \mathcal{B}$.
        Eşdeğer: $\{f \leq b\} \in \mathcal{F}$ her $b$.
        Basit fonksiyonlar artan limit.
  \item \textbf{Lebesgue integrali:} Basit $\to$ $[0,\infty]$ $\to$ $\mathbb{R}$.
        $L^1(\mu)$: $\int|f|<\infty$. Doğrusal, monoton, n.k.'ya duyarsız.
  \item \textbf{MYT:} $f_n \nearrow f \geq 0 \Rightarrow \int f_n \to \int f$.
  \item \textbf{Fatou:} $f_n \geq 0 \Rightarrow \int\liminf f_n \leq \liminf\int f_n$.
  \item \textbf{DYT:} $f_n \to f$ n.k., $|f_n| \leq g \in L^1 \Rightarrow \int f_n \to \int f$.
  \item \textbf{Fubini-Tonelli:} $\sigma$-sonlu + ($f\geq 0$ ya da $f\in L^1$)
        $\Rightarrow$ tekrarlı integral sırası serbesttir.
\end{itemize}
\end{bolumozeti}

\begin{ozyansima}
\begin{itemize}
  \item Bir $\sigma$-cebiri örnekleyin ve üç aksiyomu teker teker doğrulayın.
  \item Neden $\mu(A_1) < \infty$ koşulu olmadan yukarıdan süreklilik çöker?
        Kendi örneğinizi üretin.
  \item MYT, Fatou ve DYT'yi karşılaştırın: hangi durumlarda hangisini kullanırsınız?
  \item Fubini'nin neden $L^1$ koşulu gerektirdiğini karşı örnekle açıklayın.
\end{itemize}
\end{ozyansima}

\begin{kopru}
\textbf{Sonraki bölüm (Bölüm~2: Olasılık Uzayları).}
Bu bölümde geliştirilen ölçüm teorisi altyapısı, Bölüm~2'de olasılık ölçüsü
$\mathbb{P}$ ile hayata geçirilecek. Koşullu olasılık, bağımsızlık ve Bayes
teoremi $(\Omega,\mathcal{F},\mathbb{P})$ dili içinde yeniden formüle edilecek.

\medskip

\textbf{İleri okuma.}
Billingsley~\cite{billingsley1995} Bölüm~1--3;
Klenke~\cite{klenke2014} Bölüm~1--6;
Shiryaev~\cite{shiryaev1996} Bölüm~II §1--3.
\end{kopru}


% =====================================================================
%  ALIŞTIRMALAR
% =====================================================================

\cozumlualistirmalar

\begin{alistirma}[Al.1.1]{A}{$\sigma$-cebiri doğrulama}
$\Omega = \mathbb{R}$ ve $\mathcal{F} = \{A \subseteq \mathbb{R} : A \text{ ya da } A^c
\text{ sayılabilirdir}\}$ ailesinin bir $\sigma$-cebiri olduğunu gösterin.
\end{alistirma}
\alcozum{%
\textbf{(A)} $\mathbb{R}$ sayılamaz; dolayısıyla $\mathbb{R}^c = \emptyset$ sayılabilir.
$\mathbb{R} \in \mathcal{F}$.

\textbf{(B)} $A \in \mathcal{F}$: $A$ sayılabilirse $A^c$'nin tümleni $A$ sayılabilir,
yani $A^c \in \mathcal{F}$; $A^c$ sayılabilirse $A^c \in \mathcal{F}$ doğrudan.

\textbf{(C)} $(A_n) \subseteq \mathcal{F}$, $A = \bigcup_n A_n$.
Eğer her $A_n$ sayılabilirse $A$ sayılabilir (Teorem~0.\ref{thm:sayilabilirbir});
$A \in \mathcal{F}$.
Eğer bazı $A_{n_0}$ sayılamaz ise $A_{n_0}^c$ sayılabilir ve
$A^c = \bigcap_n A_n^c \subseteq A_{n_0}^c$ sayılabilir; $A \in \mathcal{F}$.
$\blacksquare$%
}

\begin{alistirma}[Al.1.2]{B}{Aşağıdan süreklilik uygulaması}
$(\mathbb{R}, \mathcal{B}(\mathbb{R}), \lambda)$ Lebesgue ölçüm uzayı.
$A_n = (0, 1/n)$ için $\lambda(A_n)$'i hesaplayın ve
$\lambda\!\left(\bigcap_{n=1}^\infty A_n\right)$'yi yukarıdan süreklilik
teoremi ile bulun. Doğrudan hesapla doğrulayın.
\end{alistirma}
\alcozum{%
$\lambda(A_n) = 1/n$. $A_n \searrow \emptyset$ ($\bigcap_n (0,1/n) = \emptyset$;
eğer $x > 0$ ise $x \notin (0,1/n)$ yeterince büyük $n$ için; $x \leq 0$ ise hiç içermez).
$\lambda(A_1) = 1 < \infty$; yukarıdan süreklilik:
$\lambda\!\left(\bigcap_n A_n\right) = \lim_n \lambda(A_n) = \lim_n 1/n = 0$.
Doğrudan: $\lambda(\emptyset) = 0$. Tutarlı. $\blacksquare$%
}

\begin{alistirma}[Al.1.3]{B}{MYT uygulaması}
$f_n(x) = \gostf{[0,n]}(x) \cdot e^{-x}$ ($x \geq 0$), Lebesgue ölçüsüyle.
$\int_0^\infty f_n \, d\lambda$'yı hesaplayın ve MYT kullanarak
$\int_0^\infty e^{-x} \, d\lambda$'yı bulun.
\end{alistirma}
\alcozum{%
$\int_0^\infty f_n \, d\lambda = \int_0^n e^{-x} \, dx = 1 - e^{-n}$.
$f_n \nearrow e^{-x}$ noktasal ($f_n(x) = e^{-x}$ eğer $x \leq n$, $= 0$ eğer $x > n$).
MYT: $\int_0^\infty e^{-x} \, d\lambda = \lim_n (1 - e^{-n}) = 1$. $\blacksquare$%
}

\begin{alistirma}[Al.1.4]{A}{Basit fonksiyon inşası}
$f(x) = x^2 + 1$, $x \in [0,2]$ için Teorem~\ref{thm:basityaklasim}'in
$n=2$ adımındaki $f_2$'yi açıkça yazın ve $\|f - f_2\|_\infty$'u hesaplayın.
\end{alistirma}
\alcozum{%
$n=2$: bölümleme $k/4$, $k=0,1,\ldots,7$ ve $\{f \geq 2\}$.
$f([0,2]) = [1,5]$. $f_2 = \sum_{k=0}^{7} \frac{k}{4} \gostf{\{k/4 \leq f < (k+1)/4\}} + 2\gostf{\{f \geq 2\}}$.
$f$ artan olduğundan $\{k/4 \leq f < (k+1)/4\}$ bir aralık. En büyük hata:
$f$ ile $f_2$ arasındaki fark en fazla $1/4 = 1/2^n$. $\|f-f_2\|_\infty \leq 1/4$.
$\blacksquare$%
}

\alistirmalar

\begin{alistirma}[Al.1.5]{A}{$\sigma$-cebiri işlemleri}
$\mathcal{F}_1$ ve $\mathcal{F}_2$, $\Omega$ üzerinde iki $\sigma$-cebiri.
(a)~$\mathcal{F}_1 \cap \mathcal{F}_2$'nin $\sigma$-cebiri olduğunu gösterin.
(b)~$\mathcal{F}_1 \cup \mathcal{F}_2$'nin genel olarak $\sigma$-cebiri olmayabileceğini
    bir karşı örnekle gösterin.
\end{alistirma}
\alcozum{%
(a) $A \in \mathcal{F}_1 \cap \mathcal{F}_2 \Rightarrow A \in \mathcal{F}_1$ ve
$A \in \mathcal{F}_2$; her iki $\sigma$-cebiri tümleme ve sayılabilir birleşim altında
kapalı; kesişim de kapalı.

(b) $\Omega = \{1,2,3\}$, $\mathcal{F}_1 = \{\emptyset,\{1\},\{2,3\},\Omega\}$,
$\mathcal{F}_2 = \{\emptyset,\{2\},\{1,3\},\Omega\}$.
$\{1\}, \{2\} \in \mathcal{F}_1 \cup \mathcal{F}_2$ ama
$\{1\} \cup \{2\} = \{1,2\} \notin \mathcal{F}_1 \cup \mathcal{F}_2$. $\blacksquare$%
}

\begin{alistirma}[Al.1.6]{A}{Ölçü özellikleri}
$(\Omega, \mathcal{F}, \mu)$ ölçüm uzayı, $A, B \in \mathcal{F}$ ve $\mu(A), \mu(B) < \infty$.
Gösterin ki $\mu(A \cup B) = \mu(A) + \mu(B) - \mu(A \cap B)$.
\end{alistirma}
\alcozum{%
$A \cup B = (A \setminus B) \cup (A \cap B) \cup (B \setminus A)$ ikili ayrık:
$\mu(A \cup B) = \mu(A \setminus B) + \mu(A \cap B) + \mu(B \setminus A)$.
$\mu(A) = \mu(A \setminus B) + \mu(A \cap B)$; $\mu(B) = \mu(B \setminus A) + \mu(A \cap B)$.
Toplarsak: $\mu(A) + \mu(B) = \mu(A \cup B) + \mu(A \cap B)$. $\blacksquare$%
}

\begin{alistirma}[Al.1.7]{A}{Ölçülebilir fonksiyon}
$f: \mathbb{R} \to \mathbb{R}$, $f(x) = \sin x$ fonksiyonunun
$(\mathbb{R}, \mathcal{B}(\mathbb{R}))$ üzerinde ölçülebilir olduğunu gösterin.
\end{alistirma}
\alcozum{%
$\sin: \mathbb{R} \to \mathbb{R}$ sürekli; sürekli fonksiyonlar Borel ölçülebilirdir
(açık kümenin ters görüntüsü açık, dolayısıyla Borel). $\blacksquare$%
}

\begin{alistirma}[Al.1.8]{A}{Fatou lemması --- kesin örnek}
$f_n = n \gostf{(0,1/n)}$ ($n \geq 1$), $[0,1]$ üzerinde Lebesgue ölçüsüyle.
(a)~$\liminf_n f_n$'i bulun.
(b)~$\int \liminf_n f_n$ ile $\liminf_n \int f_n$'yi karşılaştırın.
\end{alistirma}
\alcozum{%
(a) Her $x > 0$ için yeterince büyük $n$'de $x \geq 1/n$; $f_n(x) = 0$.
$f_n(0) = 0$ her zaman. $\liminf_n f_n = 0$ n.k.

(b) $\int \liminf f_n = 0$.
$\int f_n = n \cdot \lambda((0,1/n)) = n \cdot 1/n = 1$.
$\liminf_n \int f_n = 1 > 0$.
Fatou eşitsizliği $0 \leq 1$; kesin eşitsizlik. $\blacksquare$%
}

\begin{alistirma}[Al.1.9]{B}{DYT ile integral sınırı}
$f_n(x) = \dfrac{\sin(x/n)}{x/n} \cdot e^{-|x|}$ ($x \neq 0$),
$f_n(0) = e^0 = 1$, Lebesgue ölçüsüyle $\mathbb{R}$ üzerinde.
(a)~$f_n(x) \to e^{-|x|}$ noktasal olarak gösterin.
(b)~Uygun dominant fonksiyonu bulun.
(c)~$\lim_n \int_{-\infty}^\infty f_n \, d\lambda$'yı hesaplayın.
\end{alistirma}
\alcozum{%
(a) $|\sin t / t| \leq 1$ ve $\sin t / t \to 1$ ($t \to 0$); $x/n \to 0$;
$f_n(x) \to 1 \cdot e^{-|x|} = e^{-|x|}$.

(b) $|f_n(x)| \leq e^{-|x|}$ çünkü $|\sin t/t| \leq 1$.
$g(x) = e^{-|x|} \in L^1(\mathbb{R})$: $\int_{-\infty}^\infty e^{-|x|} dx = 2$.

(c) DYT: $\lim_n \int f_n = \int e^{-|x|} \, d\lambda = 2$. $\blacksquare$%
}

\begin{alistirma}[Al.1.10]{B}{Fubini ve sıra değiştirme}
$f(x,y) = e^{-xy}$, $(x,y) \in [0,\infty) \times [0,\infty)$, Lebesgue ölçüsüyle.
$\int_1^\infty \left(\int_0^\infty e^{-xy} \, dy\right) dx$'i
önce $y$'ye göre integralleyerek, ardından $x$'e göre integralleyerek hesaplayın.
Fubini'nin hangi koşulunu kontrol etmelisiniz?
\end{alistirma}
\alcozum{%
İç integral: $\int_0^\infty e^{-xy} dy = 1/x$ ($x > 0$).
Dış integral: $\int_1^\infty 1/x \, dx = \infty$.

Önce $x$: $\int_1^\infty e^{-xy} dx = e^{-y}/y$ ($y > 0$).
$\int_0^\infty e^{-y}/y \, dy$: $y \to 0^+$ yakınında $1/y$ integrallenemeyen sapma.

Fubini koşulu: $\int_1^\infty \int_0^\infty |f| \, dy \, dx = \int_1^\infty 1/x \, dx = \infty$;
$f \notin L^1([1,\infty) \times [0,\infty))$. Fubini uygulanamaz; sıra fark yaratır. $\blacksquare$%
}

\begin{alistirma}[Al.1.11]{B}{Ölçü sürekliliği --- uygulama}
$\mu$ $(\Omega,\mathcal{F})$ üzerinde bir olasılık ölçüsü. $A_n \in \mathcal{F}$
ve $\sum_{n=1}^\infty \mu(A_n) < \infty$ ise, $\mu\!\left(\limsup_n A_n\right) = 0$
olduğunu ispatlayın.
\end{alistirma}
\alcozum{%
$B_n = \bigcup_{k \geq n} A_k$ olsun; $B_n \searrow \limsup_n A_n$.
Alt-toplamsallık: $\mu(B_n) \leq \sum_{k \geq n} \mu(A_k) \to 0$
($\sum \mu(A_k) < \infty$ kuyruk toplamı).
$\mu(B_1) \leq \sum_k \mu(A_k) < \infty$; yukarıdan süreklilik:
$\mu(\limsup_n A_n) = \lim_n \mu(B_n) = 0$. $\blacksquare$

\emph{Not:} Bu, Borel-Cantelli lemmasının birinci yarısıdır (Bölüm~7).%
}

\begin{alistirma}[Al.1.12]{C}{Lebesgue--Stieltjes ölçüsü}
$F: \mathbb{R} \to \mathbb{R}$ sağdan sürekli, artan bir fonksiyon olsun.
$\mu_F((a,b]) \coloneqq F(b) - F(a)$ ile tanımlanan set fonksiyonunun
yarı-halkalar üzerinde bir ön-ölçü (pre-measure) olduğunu gösterin ve
Carathéodory genişlemesiyle $\mathbb{R}$ üzerinde bir ölçü elde edildiğini açıklayın.
$F(x) = x$ durumunda Lebesgue ölçüsü geri gelir; bu nasıl görülür?
\end{alistirma}
\alcozum{%
\begin{ipucu}{1}
$\mu_F$ için $\sigma$-toplamsallık: $(a,b] = \bigcup_n (a_n, b_n]$ ikili ayrık ise
$F(b) - F(a) = \sum_n (F(b_n) - F(a_n))$. Sağdan süreklilik bu eşitlik için gereklidir.
\end{ipucu}

$\mu_F(\emptyset) = F(a) - F(a) = 0$. Yarı-halka $\{(a,b]\}$ üzerinde $\sigma$-toplamsallık:
$(a,b]$ ikili ayrık $(a_n,b_n]$ birleşimi ise, $F$'nin sağdan sürekliliği ve Heine-Borel
gibi argümanlarla (ayrıntı için Billingsley~\cite{billingsley1995} §12 bakınız)
$F(b)-F(a) = \sum_n (F(b_n)-F(a_n))$ kanıtlanır.

Carathéodory genişlemesi bu ön-ölçüyü $\mathcal{B}(\mathbb{R})$'yi içeren bir
$\sigma$-cebire taşır. $F(x)=x$: $\mu_F((a,b]) = b-a$; bu Lebesgue dış ölçüsünü
aralıklarda tanımlar. $\blacksquare$%
}

\begin{alistirma}[Al.1.13]{C}{$L^p$ uzayları}
$(\Omega,\mathcal{F},\mu)$ ölçüm uzayı, $1 \leq p < \infty$.
$L^p(\mu) = \{f \text{ ölçülebilir} : \int |f|^p \, d\mu < \infty\}$
tanımlanıyor; $\|f\|_p = \left(\int |f|^p \, d\mu\right)^{1/p}$.

(a)~Hölder eşitsizliğini ifade edin ve ispatlayın:
$\|fg\|_1 \leq \|f\|_p \|g\|_q$ ($1/p + 1/q = 1$).

(b)~Minkowski eşitsizliğini kullanarak $\|\cdot\|_p$'nin üçgen eşitsizliğini
    sağladığını gösterin (ispat özeti yeterli).

(c)~$L^p \subseteq L^1$ olması ne zaman garantilenir? Karşı örnek veya koşul.
\end{alistirma}
\alcozum{%
(a) Young eşitsizliği: $ab \leq a^p/p + b^q/q$ ($a,b \geq 0$).
$|f(\omega)| / \|f\|_p$ ve $|g(\omega)| / \|g\|_q$'ya uygula ve integra al;
$\int |fg| / (\|f\|_p \|g\|_q) \leq 1/p + 1/q = 1$. $\|fg\|_1 \leq \|f\|_p \|g\|_q$.

(b) $\|f+g\|_p^p = \int|f+g|^p \leq \int|f+g|^{p-1}(|f|+|g|)$;
Hölder'i $|f+g|^{p-1} \in L^q$'ya uygula; $\|f+g\|_p \leq \|f\|_p + \|g\|_p$.

(c) $\mu(\Omega) < \infty$ ise Hölder ile $\|f\|_1 \leq \mu(\Omega)^{1/q} \|f\|_p < \infty$.
Sonsuz ölçüde $L^p \not\subseteq L^1$: $f(x) = (1+x)^{-1/p}$ üzerinde $[0,\infty)$,
$\int f^p = \infty$... daha dikkatli: $f(x) = x^{-1/2} \gostf{(0,1]} \in L^1 \setminus L^2$
ve $g(x) = (1+x)^{-1} \in L^2(\mathbb{R}) \setminus L^1(\mathbb{R})$. $\blacksquare$%
}

\begin{alistirma}[Al.1.14]{B}{Radon-Nikodym Türevi}
$\Omega = \{1, 2, 3, 4\}$, $\mathcal{F} = \mathcal{P}(\Omega)$.
$\mu(\{k\}) = 1/4$ (uniform), $\nu(\{1\}) = 1/2$, $\nu(\{2\}) = 1/4$, $\nu(\{3\}) = 1/8$, $\nu(\{4\}) = 1/8$.
(a) $\nu \ll \mu$ olduğunu doğrulayın. (b) Radon-Nikodym türevi $f = d\nu/d\mu$'yu bulun.
(c) $\nu(A) = \int_A f \, d\mu$ olduğunu $A = \{1,3\}$ için doğrulayın.
\end{alistirma}
\alcozum{%
(a) $\mu(A) = 0 \Rightarrow A = \emptyset \Rightarrow \nu(A) = 0$. Evet, $\nu \ll \mu$.\\
(b) RN türevi: $f(\omega) = \nu(\{\omega\})/\mu(\{\omega\})$ (ayrık durumda).
$f(1) = (1/2)/(1/4) = 2$; $f(2) = 1$; $f(3) = 1/2$; $f(4) = 1/2$.\\
(c) $\nu(\{1,3\}) = 1/2 + 1/8 = 5/8$.
$\int_{\{1,3\}} f \, d\mu = f(1)\mu(\{1\}) + f(3)\mu(\{3\}) = 2\cdot(1/4) + (1/2)\cdot(1/4) = 1/2 + 1/8 = 5/8$. \checkmark}

\begin{alistirma}[Al.1.15]{B}{$L^p$ Norm Hesabı}
$f(x) = e^{-x}$, $x \geq 0$; Lebesgue ölçüsü.
(a) $\|f\|_p = (\int_0^\infty e^{-px} dx)^{1/p}$'yi hesaplayın.
(b) $p \to \infty$ limitinde $\|f\|_p$'nin ne olduğunu yorumlayın.
(c) Hölder eşitsizliğiyle $\int_0^\infty e^{-x}\cos x \, dx \leq \|f\|_2\|\cos\|_{L^2([0,\infty))}$
    üst sınırını verin (dikkat: $\|\cos\|_{L^2([0,\infty))} = \infty$; bu ne anlama gelir?).
\end{alistirma}
\alcozum{%
(a) $\|f\|_p^p = \int_0^\infty e^{-px} dx = 1/p$. $\|f\|_p = p^{-1/p}$.\\
(b) $\|f\|_\infty = \sup_{x \geq 0} |e^{-x}| = e^0 = 1$. $\lim_{p\to\infty} p^{-1/p} = 1$ (çünkü $\ln(p^{-1/p}) = -\ln p/p \to 0$). $L^\infty$ normu esansiyel supremum.\\
(c) $\cos \notin L^2([0,\infty))$ çünkü $\int_0^\infty \cos^2 x \, dx = \infty$ (yarı-periyot toplamı ıraksak). Hölder uygulanamaz — $g \in L^q$ gerektirir. Bu, Hölder'in koşulsuz uygulanamayacağını gösterir.}
